\documentclass[10pt,a4paper]{article}

\usepackage[utf8]{inputenc}
\usepackage[T1]{fontenc}
\usepackage[french]{babel}
\usepackage{amsmath,amssymb,amsthm}
\usepackage{geometry}
\usepackage{booktabs}
\usepackage{enumitem}
\usepackage{tikz}
\usepackage{hyperref}

\geometry{hmargin=1.5cm, top=1cm, bottom=1.2cm, footskip=0.5cm}
\setlength{\parskip}{3pt}
\setlength{\parindent}{0pt}
\setlist{itemsep=1pt,topsep=2pt}

\newcommand{\A}{\mathcal{A}}
\newcommand{\R}{\mathcal{R}}
\newcommand{\N}{\mathcal{N}}

\title{Projet SDP -- Apprentissage de modèles MR-Sort}
\author{Côme-Alexis Puech \and Hugo Vigna \and Salomé Fonvielle}
\date{Octobre 2026}

\begin{document}

\maketitle

\section*{Préliminaires : admission d'étudiants}

% ------------------------------------------------------------------
\subsection*{Question i}

Somme des poids des cours validés ($\geq 12$), avec $w=(0.25,0.15,0.2,0.4)$, $\lambda=0.66$ :
Alice $0.15+0.2+0.4=0.75$ ($\A$), Bob $0.25+0.15+0.2=0.60$ ($\R$),
Charlie $0.15+0.4=0.55$ ($\R$), David $1$ ($\A$).
Bob ne peut pas améliorer sa situation : avec 12 en physique, ce cours est déjà validé,
une meilleure note n'y change rien (seul H, à 10, pourrait le faire admettre).

% ------------------------------------------------------------------
\subsection*{Question ii}

Les 4 cours doivent jouer des rôles symétriques (3 cours quelconques suffisent), d'où
$w_M=w_P=w_L=w_H=\frac14$ et $\lambda\in\left]\frac12,\frac34\right]$, par exemple $\lambda=0.75$ :
3 cours pèsent $\frac34\geq\lambda$, 2 cours pèsent $\frac12<\lambda$.

% ------------------------------------------------------------------
\subsection*{Question iii}

On liste les coalitions gagnantes (somme des poids $\geq \lambda$) dans les deux cas.

\begin{center}
\begin{tabular}{lcc}
\toprule
Coalition & $w = (0.49, 0.49, 0.02)$, $\lambda=0.5$ & $w = (\frac13,\frac13,\frac13)$, $\lambda=0.5$ \\
\midrule
$\{1\}$, $\{2\}$, $\{3\}$ & $0.49,\ 0.49,\ 0.02$ : perdantes & $0.33$ : perdantes \\
$\{1,2\}$ & $0.98$ : gagnante & $0.67$ : gagnante \\
$\{1,3\}$, $\{2,3\}$ & $0.51$ : gagnantes & $0.67$ : gagnantes \\
$\{1,2,3\}$ & $1$ : gagnante & $1$ : gagnante \\
\bottomrule
\end{tabular}
\end{center}

Dans les deux cas, un étudiant est admis si et seulement s'il valide
\emph{au moins 2 cours sur 3}, quels qu'ils soient : les deux jeux de paramètres
définissent exactement la même règle de décision.

\textbf{Conclusion.} $(w,\lambda)$ n'est pas unique : seul compte l'ensemble des coalitions gagnantes.
Les poids ne mesurent pas l'importance (le cours 3, de poids $0.02$, est aussi décisif que les autres) ;
un modèle appris s'évalue donc sur ses décisions, pas sur ses paramètres.

% ------------------------------------------------------------------
\subsection*{Question iv}

Comme $b$ est connu, on sait pour chaque objet $x$ quels critères il valide : $C(x)=\{i\in\N : x_i\geq b_i\}$.
Son score $s(x)=\sum_{i\in C(x)} w_i$ est donc linéaire en $w$.
\textbf{Variables :} $w_i\geq 0$, $\lambda\in[0,1]$ et $\alpha$, la marge.
\begin{align*}
\max\ & \alpha \\
\text{s.c.}\ & s(x) \geq \lambda + \alpha && \forall x\in A^* \\
& s(x) \leq \lambda - \alpha && \forall x\in R^* \\
& \textstyle\sum_{i\in\N} w_i = 1
\end{align*}
$\alpha$ est l'écart minimal entre les scores et le seuil $\lambda$ : on cherche le modèle qui sépare le plus nettement
acceptés et rejetés. Si $\alpha^*>0$, tous les objets sont bien classés ; sinon, aucun modèle ne restitue parfaitement les données.

\begin{center}
\begin{tikzpicture}[x=26cm,y=1cm,font=\small]
  \fill[gray!15] (0.60,-0.15) rectangle (0.75,0.9);
  \draw[->] (0.48,0) -- (1.03,0) node[right] {$s(x)$};
  \foreach \t in {0.5,0.6,0.7,0.8,0.9,1.0} \draw (\t,0.06) -- (\t,-0.06) node[below] {\footnotesize\t};
  \draw[dashed,thick] (0.675,-0.15) -- (0.675,1.15) node[above] {$\lambda$};
  \draw[<->] (0.60,0.7) -- node[above] {$\alpha$} (0.675,0.7);
  \draw[<->] (0.675,0.7) -- node[above] {$\alpha$} (0.75,0.7);
  \foreach \v/\n in {0.55/Charlie,0.60/Bob} \fill[red!75!black] (\v,0) circle (2.2pt) node[above=2pt,text=black] {\n};
  \foreach \v/\n in {0.75/Alice,1.0/David} \fill[green!50!black] (\v,0) circle (2.2pt) node[above=2pt,text=black] {\n};
\end{tikzpicture}

\footnotesize Exemple de la question i : rejetés en rouge, acceptés en vert.
\end{center}

% ------------------------------------------------------------------
\subsection*{Question v}

% TODO : contraintes définissant delta_i(x).

% ------------------------------------------------------------------
\subsection*{Question vi}

% TODO : linéarisation de w_i(x) et condition d'admission.

% ------------------------------------------------------------------
\subsection*{Question vii}

% TODO : MILP apprenant simultanément b, w et lambda.

% ------------------------------------------------------------------
\subsection*{Question viii -- Résolution avec Gurobi}

% TODO

% ------------------------------------------------------------------
\subsection*{Question viii bis -- Généralisation à $p$ catégories}

% TODO

\end{document}
